\documentclass[12pt, a4paper]{article}
\usepackage[utf8]{inputenc}
\usepackage[T1]{fontenc}
\usepackage{geometry}      
\geometry{a4paper, margin=1in}
\usepackage{titlesec}      
\usepackage{fancyhdr}    
\usepackage{graphicx}       
\usepackage{hyperref}       
\usepackage{array, booktabs}
\usepackage{tabularx}       
\usepackage{enumitem}       
\usepackage{xcolor}       
\usepackage[most]{tcolorbox}
\titlespacing*{\section}{0pt}{3.5ex plus 1ex minus .2ex}{2.3ex plus .2ex}
\titlespacing*{\subsection}{0pt}{3.25ex plus 1ex minus .2ex}{1.5ex plus .2ex}
\pagestyle{fancy}
\fancyhf{} 
\fancyhead[L]{Write title here} 
\fancyhead[R]{\thepage}  
\renewcommand{\headrulewidth}{0.4pt}

\hypersetup{
    colorlinks=true,
    linkcolor=black,
    urlcolor=blue!70!black,
    pdftitle={Lab Report},
}

\begin{document}

% ==========================================
% TITLE PAGE
% ==========================================
\begin{titlepage}
    \centering
    
    % University Header
    {\LARGE \textbf{{SIKSHA `O' ANUSANDHAN}} \par}
    {\Large \textbf{(DEEMED TO BE UNIVERSITY)} \par}
    \vspace{1.5cm}
    
    % Batch and Session Row
    \noindent
    {\large \textbf{{Admission Batch:}} \underline{\hspace{4cm}}} \hfill {\large \textbf{{Session:}} \underline{\hspace{4cm}}}
    \vspace{1.5cm}
    
    % Document Title
    {\large \textbf{Laboratory Record} \par}
    {\LARGE \textbf{Machine Learning Projects with Python} \par}
    {\LARGE \textbf{(CSE 4192)} \par}
    \vspace{1cm}
    
        \begin{flushleft}
            \hspace{1cm} {\large \textbf{Submitted by}} \par    
            \vspace{0.25cm}        
            \hspace{4cm} 1. Name: Student 1 (Regd. No.: 01) \par
            \vspace{0.25cm}        
            \hspace{4cm} 2. Name: Student 2 (Regd. No.: 02) \par
            \vspace{0.25cm}        
            \hspace{4cm} 3. Name: Student 3 (Regd. No.: 03) \par
            \vspace{0.25cm}
            \hspace{4cm} 4. Name: Student 4 (Regd. No.: 04) \par
        \end{flushleft}
   
	\begin{center}
	\includegraphics[width=3cm]{logo.png}
    \end{center}
    \vspace{1.5cm}
    \vspace{1.5cm}
    {\Large \textbf{Centre for Artificcial Intelligence \& Machine Learning} \par}
    \vspace{0.3cm}
    {\Large \textbf{Faculty of Engineering \& Technology (ITER)} \par}
    \vspace{0.3cm}
    {\large \textbf{Jagamohan Nagar, Jagamara, Bhubaneswar, Odisha - 751030} \par}
\end{titlepage}

    
    
   


% ==========================================
% TABLE OF CONTENTS
% ==========================================
% \newpage
% \tableofcontents 
% \newpage

% ==========================================
% MAIN CONTENT
% ==========================================
\section{Introduction}
\subsection{Problem Statement}
\textit{Write your problem statement here. Adding text ensures the section renders fully in the PDF.}

\subsection{Aim}
\textit{State the primary aim of your project.}

\subsection{Objectives}
\textit{List the specific objectives.}

\section{Dataset Description}
\subsection{Dataset Source}
\textit{Provide content here.}

\subsection{Dataset Features}
\textit{Provide content here.}

\subsection{Target Variable}
\textit{Provide content here.}

\subsection{Dataset Statistics}
\textit{Provide content here.}

\section{Exploratory Data Analysis}
\subsection{Dataset Structure}
\textit{Provide content here.}

\subsection{Descriptive Statistics}
\textit{Provide content here.}

\subsection{Missing Value Analysis}
\textit{Provide content here.}

\subsection{Class Distribution}
\textit{Provide content here.}

\subsection{Univariate Analysis}
\textit{Provide content here.}

\subsection{Bivariate Analysis}
\textit{Provide content here.}

\subsection{Correlation Analysis}
\textit{Provide content here.}

\subsection{Outlier Analysis}
\textit{Provide content here.}

\section{Data Preprocessing}
\subsection{Data Cleaning}
\textit{Provide content here.}

\subsection{Missing Value Handling}
\textit{Provide content here.}

\subsection{Outlier Handling}
\textit{Provide content here.}

\subsection{Feature Engineering}
\textit{Provide content here.}

\subsection{Feature Selection}
\textit{Provide content here.}

\subsection{Train-Test Split}
\textit{Provide content here.}

\subsection{Feature Scaling}
\textit{Provide content here.}

\section{Methodology}
\subsection{Overall Project Workflow}
\textit{Provide content here.}

\subsection{Model Development Strategy}
\textit{Provide content here.}

\subsection{Baseline Models}
\textit{Provide content here.}

\subsection{Multilayer Perceptron Architecture}
\textit{Provide content here.}

\section{Model Development}
\subsection{Logistic Regression}
\textit{Provide content here.}

\subsection{K-Nearest Neighbors}
\textit{Provide content here.}

\subsection{Support Vector Machine}
\textit{Provide content here.}

\subsection{Random Forest}
\textit{Provide content here.}

\subsection{Multilayer Perceptron}
\textit{Provide content here.}

\section{Model Evaluation}
\subsection{Accuracy}
\textit{Provide content here.}

\subsection{Precision}
\textit{Provide content here.}

\subsection{Recall}
\textit{Provide content here.}

\subsection{F1-Score}
\textit{Provide content here.}

\subsection{Confusion Matrix}
\textit{Provide content here.}

\subsection{ROC-AUC}
\textit{Provide content here.}

\subsection{ROC Curve}
\textit{Provide content here.}

\subsection{Precision-Recall Curve}
\textit{Provide content here.}

\section{Model Comparison and Selection}
\subsection{Performance Comparison}
\textit{Provide content here.}

\subsection{Model Selection Criteria}
\textit{Provide content here.}

\subsection{Selected Model}
\textit{Provide content here.}

\section{Hyperparameter Tuning}
\subsection{Hyperparameters Considered}
\textit{Provide content here.}

\subsection{Tuning Method}
\textit{Provide content here.}

\subsection{Search Space}
\textit{Provide content here.}

\subsection{Tuning Results}
\textit{Provide content here.}

\subsection{Optimized Architecture}
\textit{Provide content here.}

\section{Final Model Training}
\textit{Provide content here.}

\section{Final Model Evaluation}
\textit{Provide content here.}

\section{Model Deployment and Testing}
\subsection{Deployment Approach}
\textit{Provide content here.}

\subsection{Deployment Tool}
\textit{Provide content here.}

\subsection{User Interface}
\textit{Provide content here.}

\subsection{Prediction Workflow}
\textit{Provide content here.}

\subsection{Deployment Testing}
\textit{Provide content here.}

\section{Results and Discussion}
Write your results here.

\section{Limitations}
\textit{Provide content here.}

\section{Conclusion}
\textit{Provide content here.}

\section{Future Scope}
\textit{Provide content here.}


\vspace{2cm}

\noindent \textit{To cite the references, you can use below given method.}


\cite{ref1} %to cite a reference
\cite{ref2}
\cite{ref3}
\begin{thebibliography}{99}
\bibitem{ref1} 
I. Goodfellow, Y. Bengio, and A. Courville, \textit{Deep Learning}, MIT Press, 2016.

\bibitem{ref2} 
J. Smith and A. Doe, ``A Novel Approach to Gamified Mathematics Education,'' \textit{Journal of Educational Computing}, vol. 45, no. 2, pp. 112-130, 2023.

\bibitem{ref3} 
National Council of Educational Research and Training (NCERT), \textit{Mathematics Textbook for Class IX}, New Delhi: NCERT, 2022.

\end{thebibliography}


\appendix 
\section{Source Code}
\textit{Provide content here.}

\section{Additional Results}
\textit{Provide content here.}

\section{Deployment Screenshots}
\textit{Provide content here.}

\section{Contributions of Group Members}

\noindent
\begin{tabularx}{\textwidth}{|c|X|c|X|c|}
\hline
\textbf{Sl.} & \textbf{Group Member} & \textbf{Regd. No.} & \textbf{Role / Responsibility} & \textbf{Contribution (\%)} \\
\hline
1 & & & & \\
\hline
2 & & & & \\
\hline
3 & & & & \\
\hline
4 & & & & \\
\hline
\end{tabularx}

\vspace{0.5cm}
\noindent \textit{Students should indicate individual responsibility for areas such as:}
\begin{itemize}[noitemsep, topsep=5pt]
    \item Dataset collection and understanding 
    \item Exploratory Data Analysis 
    \item Data preprocessing 
    \item Feature engineering 
    \item Baseline model development 
    \item MLP development 
    \item Hyperparameter tuning 
    \item Model evaluation and comparison 
    \item Deployment 
    \item Report preparation 
\end{itemize}

\end{document}
